\documentclass[11pt]{letter}
\usepackage[margin=2.2cm]{geometry}
\usepackage{mathptmx}
\usepackage[hidelinks]{hyperref}
\signature{P. N. Bhargav Teja, Lanka Sree Chathurya, K. Arun Reddy and Ragavan K\\
School of Computer Science and Engineering, Vellore Institute of Technology, Vellore, India\\Corresponding author: Ragavan K, [ADD E-MAIL]}
\address{School of Computer Science and Engineering\\Vellore Institute of Technology\\Vellore, Tamil Nadu 632014, India}
\begin{document}
\begin{letter}{The Editor-in-Chief\\IEEE Access}
\opening{Dear Editor,}

We submit our manuscript, \emph{``Excitation-Coherent Telemetry Analysis for
Predictive OBD-II Connector Wear Detection: A Physics-Informed Study on Real
In-Vehicle Network Data,''} for consideration as a regular article in
\emph{IEEE Access}.

The SAE~J1962 diagnostic connector now serves as the permanent mount and
power source of telematics, insurance and fleet devices, although it was
designed for occasional workshop use. Vibration-induced fretting of its
contacts causes intermittent data loss and device brown-outs that are
routinely blamed on the device or the vehicle. To our knowledge, this is
the first study of whether the wear of the diagnostic connector itself can
be detected, attributed and predicted from the telemetry that the attached
device already records, without added hardware.

The main contributions are:
\begin{itemize}
\item a physics-informed model linking contact wear to the observable
  behaviour of three classes of diagnostic tool, with four realistic
  confounders;
\item ECTA, a detection and attribution method built around a test of
  whether data losses co-vary with the vehicle's own measured engine and road
  excitation;
\item a first-principles observability law that predicts, without any
  classifier, which wear stages each class of tool can detect (Spearman
  0.93 with the measured performance);
\item an evaluation on five public real-vehicle datasets (nine vehicles in
  detail and a 381-vehicle fleet) with strictly vehicle-disjoint splits and
  eleven baseline and ablation models. On unseen vehicles a full-bus tool
  detects wear with AUROC 0.880 (moderate 0.891, severe 0.995) and confirms
  moderate wear within two minutes of driving, whereas deep sequence
  baselines reach 0.646.
\end{itemize}

We have taken care to report the limits of the approach as clearly as its
successes: the connector faults are modelled on top of real telemetry (no
public dataset of naturally worn J1962 connectors exists), the unmeasured
constants of the model are stress-tested by deliberate mis-specification
and by a structurally different fault process (on which a full-bus tool
still reaches AUROC 0.856),
and the negative findings (incipient wear, moderate wear for low-rate
tools, and months-ahead remaining-life estimation) are reported in full.
A Supplementary Material document reports the complete results, paired
significance tests against every baseline, hyperparameter and window-length
sensitivity, and false-alarm rates per hour of real driving. All code is
available, and every reported number can be reproduced from it.

The work spans intelligent transportation, vehicular electronics,
fault diagnosis and prognostics and health management, which we believe
fits the broad, multidisciplinary scope of \emph{IEEE Access}.

We confirm that this manuscript is original, has not been published
previously, and is not under consideration for publication elsewhere. All
authors have approved the manuscript and agree with its submission. The
method is related to Indian Patent Application 202641074361, filed by the
authors, which is disclosed in the manuscript. The authors declare no other
competing interests.

\closing{Sincerely,}
\end{letter}
\end{document}
